\chapter{Algemene Informatie}
\label{chap:algemeneinformatie}

Dit werkboek is bedoeld voor 1e-jaars studenten aan de opleiding ICT en de Ad ITSM. Voor deze module gelden geen instroomeisen.

\section*{Leeruitkomsten}
Binnen het vak Computer Architecture maakt je kennis met de algemene architectuur van een computersysteem. Uiteindelijk kun je de algemene werking van een processor benoemen en hoe deze aangestuurd wordt benoemen. Je begrijpt de werking van de verschillende typen computergeheugen en kan hier berekeningen op uitvoeren. Daarnaast begrijpt je assembly (voor de MOS6502) en opcodes en kan zelf eenvoudige programma’s hiermee realiseren. 


\section{Introductie tot het vak}
Deze introductie dient om de lezer te informeren over de bedoeling van het vak, de opzet van het materiaal en over de manier waarop dit materiaal bestudeerd kan worden. Deze introductie bevat nog geen echte studiestof, maar meer praktische en studiegerichte informatie.

\subsection{Motivering voor het vak}
Tijdens het eerste jaar van de opleidingen ICT en de Ad ITSM wordt de basis van programmeren onderwezen. Dit wordt onder andere gedaan door studenten simpele computerprogramma’s te laten schrijven, welke dienen als middel om het schrijven van code en functies aan te leren. Om een programma te kunnen programmeren, is het noodzakelijk om ook kennis te hebben van het platform waar dit programma op zal werken. Dit vak voorziet de student daarom van de basisbeginselen binnen de wereld van computers.

% \subsection{Plaatsing van het vak}
% Dit vak is een verplicht deel van de bacheloropleiding ICT. Omdat de basiskennis belangrijk is voor beide specialisaties, is deze vroeg in het eerste jaar opgenomen als vak in het gezamenlijk deel.

% \section{De inhoud van het vak}

\subsection{Leerdoelen}

Na het bestuderen van het vak wordt verwacht dat je:
\begin{itemize}
    \item[--] de algemene architectuur van een computersysteem begrijpt
    \item[--] de algemene werking van een processor en hoe deze aangestuurd wordt kent
    \item[--] de werking van de verschillende typen computergeheugen kent
    \item[--] simpele programma's kunt maken voor de MOS 6502 processor.
    \item[--] basisbegrip hebt van digitale logica en hier kleine schakeling kan herkennen en maken.
\end{itemize}

\subsection{Opbouw van het vak}
Het werkboek is verdeeld in 4 verschillende hoofdstukken, die overeen komen met de 4 gegeven lessen. Elk hoofdstuk is een samenhangend geheel van leerdoelen en lesstof, die nodig zijn om bij afsluiting van het vak aan de gestelde eisen van tentaminering te kunnen voldoen.

\section{Aanwijzingen voor het bestuderen van het vak}
Dit vak bestaat uit een aantal onderdelen, waarvan een aantal digitaal gratis verkrijgbaar zijn via de online dienst van O’Reilly (waarover hieronder meer uitleg) of worden beschikbaar gesteld via de digitale systemen van de school. 

Binnen het vak gebruiken we het volgend tekstboek: 
\textit{Modern Computer Architecture and Organization} van Jim Ledin \cite{Ledin} (ISBN: 978-1-80323-451-9, verkrijgbaar via \href{https://learning.oreilly.com/}{O'Reilly Learning})

Het werkboek heeft de functie als leidraad te dienen tijdens het bestuderen van de lesstof. Ook staat hier aanvullende lesstof en worden opgaven gegeven. 
% Wanneer we refereren aan het tekstboek hierboven, gebruiken we de achternaam van de eerste auteur. Bijvoorbeeld, wanneer er verwezen wordt naar een passage uit het tekstboek van Jim Ledin, dan staat er
% in de tekst ``Lees paragraaf ... uit \textbf{Ledin}''.

\studeeraanwijzing{We hanteren dit symbool, om aan te geven dat we verwijzen naar het tekstboek van \textbf{Ledin}}

\subsection{Boeken via O'Reilly}
De bibliotheek van NHL Stenden heeft een overeenkomst met de uitgever O’Reilly, waarbij studenten en medewerkers van de hogeschool onbeperkte toegang hebben tot de online bibliotheek van O’Reilly. Via dit systeem kunnen de boeken die onderdeel zijn van dit vak worden verkregen. Het is natuurlijk ook mogelijk om de boeken fysiek aan te schaffen, wanneer dit de voorkeur heeft, maar dit is slechts optioneel om extra studiekosten zo laag mogelijk te houden.

Het platform is voor studenten van de NHL Stenden de eerste keer te bereiken via
\url{https://www.oreilly.com/library/view/temporary-access/}. 

Daarna kun je bij ``Select your institution'', ``NHL Stenden University of Applied Sciences'' kizen. Vervolgens word je gevraagd om je NHL Stenden e-mailadres in te voeren. Nadat je dit gedaan hebt, word je doorgelaten tot de digitale bibliotheek.

Merk hierbij op dat je nooit gevraagd is om een wachtwoord in te voeren. Dit is maar met moeite een feature te noemen in plaats van een bug. Zodra je in de digitale bibliotheek bent, kun je het beste uitloggen en naar de normale inlogpagina gaan, te bereiken via deze link: \url{https://www.oreilly.com/member/login/}. Hier kun je klikken op ``Forgot password'' en jezelf een password reset link sturen. Reset je wachtwoord en je hebt je eigen persoonlijke account.

\section*{Over dit werboek}
Dit werboek is geschreven in \LaTeX\ (spreek uit: lateg) met behulp van \href{https://neovim.io/}{Neovim} en \href{https://www.overleaf.com}{Overleaf}. 

